\def\topic{Historical Variants of Iota and Family Evaluation}
\ifdefined\MyHeaderIncluded
\else
\def\qwertyflag{true}

\ifdefined\MyHeaderIncluded

\else
  \newcommand{\foobar}{FooBar}
  \def\MyHeaderIncluded{true}
% if is till the end of the file

\documentclass[10pt,a4paper]{article}
\usepackage[utf8]{inputenc}
%\usepackage[german]{babel}
%\usepackage[cyrillic]{babel}
\usepackage[T1]{fontenc}
%\usepackage[T2A]{fontenc}
%\usepackage[english, russian]{babel}
\usepackage{amsmath}
\usepackage{amsfonts}
\usepackage{amssymb}
\usepackage{amsthm}
\usepackage{mathtools} %coloneqq

\usepackage{bussproofs} %

%\renewcommand{\thesubsection}{\arabic{subsection}}
%\newtheorem{thm}{Theorem}[section]
%\renewcommand{\thethm}{\arabic{subsection}.\arabic{thm}}

\usepackage{thmtools}
\usepackage{hyperref}
\hypersetup{
    colorlinks=true,
    linkcolor=blue,
    filecolor=magenta,      
    urlcolor=cyan,
    pdftitle={Overleaf Example},
    %pdfpagemode=FullScreen,
    }
%\usepackage{stix-type1}%\turnediota
\theoremstyle{definition}

\newtheorem{constructor}{Constructor}[section]
\newtheorem{iaxiom}{Axiom}
\newtheorem{axiom}{Axiom}
\newtheorem*{axiom*}{Axiom}
\newtheorem{definition}{Definition}[section]
\newtheorem*{definition*}{Definition}
\newtheorem{remark}{Remark}[section]
\newtheorem*{remark*}{Remark}
\newtheorem{recall}{Recall}[section]
\newtheorem*{recall*}{Recall}

\newtheorem{lemma}{Lemma}[section]
\newtheorem{theorem}{Theorem}[section]
\newtheorem*{theorem*}{Theorem}
\newtheorem{note}{Note}[section]
\usepackage{etoolbox}

%\newcommand*{\@Dashv}[2]{%
%  \reflectbox{$\m@th#1#2$}%
%}
\newcommand*{\separation}[3]{% 1)binder, 2)set, 3)condition
  \{#1\in #2: #3\}
}

% https://www.freecodecamp.org/news/simple-regex-tricks-for-beginners-3acb3fa257cb/
%  \\mathcal\{T\} |---> \topology
%  \\mathcal\{([^\}]*)\}   |--->  $1
\newcommand{\AND}{\ \ \&\ \ }
\newcommand{\any}{\mathrm{\ any}}
\newcommand{\borel}{\mathcal{B}}
\newcommand{\can}{\mathrm{can}}
\newcommand{\canonize}{\mathrm{canonize}}
\newcommand{\closure}{\mbox{cl}}
\newcommand{\complex}{\mathbb{C}}
\newcommand{\converges}{\longrightarrow}
\newcommand{\countable}{\mbox{\ countable}}
\newcommand{\ldefined}{\mbox{\ defined}}
\newcommand{\dom}{\mbox{dom}}
\newcommand{\Dom}{\mbox{Dom}}
\newcommand{\Done}{\noindent \textbf{Done.}\\}
\newcommand{\lexists}{\mbox{\ exists}}
\newcommand{\lforall}{\mbox{\ forall}}
\newcommand{\lunique}{\mbox{\ unique}}
\newcommand{\family}{\mathrm{\ family}}
\newcommand{\Fm}{\mathrm{Fm}}
\newcommand{\FuncSymb}{\mathrm{FuncSymb}}
\newcommand{\function}{\mathrm{\ function}}
\newcommand{\functional}{\mbox{\ functional}}
\newcommand{\IFF}{\ \Longleftrightarrow\ }
\renewcommand{\iff}{\ \longleftrightarrow\ }
\newcommand{\IMPLIES}{\ \Longrightarrow\ }
\renewcommand{\implies}{\ \longrightarrow\ }
\newcommand{\Inf}{\mathrm{Inf}}
\newcommand{\injective}{\mbox{\ injective}}
\newcommand{\injection}{\mbox{\ injection}}
\newcommand{\interior}{\mbox{int}}
\newcommand{\intersection}{\bigcap}
\newcommand{\isempty}{\mathrm{\ empty}}
\newcommand{\LB}{\mathrm{LB}}
\newcommand{\nghbset}{\mathcal{N}}
\newcommand{\nonempty}{\mathrm{\ nonempty}}
\newcommand{\Ord}{\mathrm{Ord}}
\newcommand{\power}{\mathcal{P}}
\newcommand{\pr}{\mathrm{pr}}
\newcommand{\fst}{\mathrm{fst}}
\newcommand{\snd}{\mathrm{snd}}
\newcommand{\PredSymb}{\mathrm{PredSymb}}
\newcommand{\Proof}{\noindent \textbf{Proof.}\\}
\newcommand{\ran}{\mathrm{ran}}
\newcommand{\Ran}{\mathrm{Ran}}
\newcommand{\real}{\mathbb{R}}
\newcommand{\relation}{\mathrm{\ relation}}
\newcommand{\set}{\mathrm{\ set}}
\newcommand{\singleton}{\mathrm{\ singleton}}
\newcommand{\semantics}{\mathrm{semantics}}
\newcommand{\semialgebra}{\mathcal{S}}
\newcommand{\Sup}{\mathrm{Sup}}
\newcommand{\surjective}{\mbox{\ surjective\ on\ }}
\newcommand{\Term}{\mathrm{Term}}
\newcommand{\Type}{\mathrm{Type}}
\newcommand{\topology}{\mathcal{T}}
\newcommand{\total}{\mbox{\ total\ on\ }}
\newcommand{\transitive}{\mathrm{\ transitive}}
\newcommand{\union}{\bigcup}
\newcommand{\lTrue}{\mathrm{True}}
\newcommand{\lFalse}{\mathrm{False}}
\newcommand{\lUncertain}{\mathrm{Uncertain}}
\newcommand{\lif}{\mathrm{if}}
\newcommand{\lthen}{\mathrm{then}}
\newcommand{\lelse}{\mathrm{else}}
\newcommand{\lexception}{\mathrm{exception}}
\newcommand{\lundef}{\mathrm{undef}}
\newcommand{\univocal}{\mathrm{\ univocal}}
\newcommand{\univocalat}{\mathrm{\ univocal\ at\ }}
\newcommand{\UB}{\mathrm{UB}}
\newcommand{\Un}{\mathcal{U}}
\newcommand{\Var}{\mathrm{Var}}

\newcommand{\notvarphi}{\phi}
\renewcommand{\phi}{\varphi}

\newcommand{\comment}[1]{\quad\mbox{(#1)}}
% local commands
%\preto\section{\undef\blipo} %etoolbox
\preto\subsection{\undef\blipo} %etoolbox
% https://tex.stackexchange.com/questions/105684/limit-macro-newcommand-scope-to-section
\usepackage{graphicx}
\usepackage{authblk}

\setlength\parindent{0pt} % whole doc of \noindent
\usepackage{soul}
%\usepackage{mathgl}

\fi % guard
\ifdefined\topic
\else
\def\topic{No Titel}
\fi
\title{\topic}
\author[1]{Georgy Dunaev}
\begin{document}
\maketitle
\fi
\section{[OUTDATED] Variants of Iota Definitions}
These notes investigate two related questions: how a relation can represent
a family with class-valued members, and how a definite descriptor can signal
that no unique set witness exists. The subsections record successive local
conventions, rather than cumulative definitions of one evaluation operation.
In particular, union evaluation, protected-family evaluation, and the
intersection-based iota descriptor have different default behavior.

The symbol $!$ means at-most-one uniqueness. Class comprehension collects
set witnesses, even when the surrounding variables range over classes.
In the final construction, $\intersection\emptyclass=V$ supplies a
proper-class default, so that sethood detects successful description.
That convention belongs to this construction and is independent of the
partial-logic treatment elsewhere in the book.
\subsection{First Attempt: Evaluating Relations and Families}
The initial proposal compares several encodings. Its central difficulty is
preserving an index whose intended value is empty: a representation made
only of members of that value loses the index when there are no members.
\begin{definition}[value of relation/function/protected family at $x$]
\[
F(x)\ceq \union\cc{b}{(x,b)\in F}
\]
\end{definition}
\begin{definition}[value of a simple family at $x$]
\[
F(x) \ceq \cc{b}{(x,b)\in F}
\]
\end{definition}
This is the fibre at $x$, whereas $\ran(F)$ collects second components
at all indices. With the ordinary relational domain, every represented
fibre is nonempty. Declaring every set to be an index permits empty fibres,
but then the domain no longer records which indices were intended.
The following encoding introduces a separate record for an empty value.
This distinction also matters in the newer Obelisk catalog:
\texttt{DefValueAtArg} returns $V$ on an absent fibre, whereas the
unconditional union convention in this subsection returns $\emptyclass$.
The two conventions agree when a record exists. A theorem transferring
between them must therefore state its domain hypothesis.
\begin{definition}[value of a protected family at $x$]
\[
F(x)_p \ceq \cc{b}{(x,\{b\})\in F}
= \union\cc{s}{(x,s)\in F}
\]
For a protected family, the second components are empty sets or singletons,
so the two displayed expressions agree. The union expression motivates a
common evaluation notation for the encodings under consideration.
\end{definition}


\begin{definition}[Domain of relation/function/family/protected family]
\[
\Dom(F)\ceq \cc{x}{\exists s,\ (x,s)\in F}=:\dom(F)
\]
\end{definition}

\begin{definition}[Range of relation]
The relational range collects the stored second components. It can differ
from the collection of values obtained by evaluating the family.
\[
\ran(R)\ceq \cc{y}{\exists x,\ (x,y)\in R}
\]
\end{definition}

\begin{definition}[in range of relation/function/family]
This predicate describes evaluated values, including proper classes.
It does not collect proper classes as members of an object-theoretic class.
\[
(K\mbox{ in classrange of }F) \IFF \exists x,\  (x \in \Dom(F)) \land (F(x)=K)
\]
\end{definition}
\begin{definition}[in range of protected family]
This predicate describes evaluated values, including proper classes.
It does not collect proper classes as members of an object-theoretic class.
\[
(K\mbox{ in classrange of protected }F) \IFF \exists x,\ (x \in \Dom(F)) \land (F(x)_p=K)
\]
\end{definition}


\begin{definition}[(Set--)Range of function/family. "F--Ran","Ran", "big range"]
%\[
%\Ran(F)\ceq \cc{y}{\exists x,\  (x \in \Dom(F)) \land (F(x)=y)}
%\]
\[
\Ran(F)\ceq \cc{y}{ y\mbox{ in classrange of } F}
\]
\end{definition}

\begin{definition}[function]
A function is a relation which is univocal.
\end{definition}
\begin{definition}[protected family]
A protected family is a relation whose second components are subsets of
singletons, hence either empty sets or singletons.
\end{definition}


%A relation is a function when it is univocal.
%A relation is a family when it's relation range consists only of subsets of singletones.

\begin{definition}[r/f/F is defined at $x$]
\[
(F \ldefined\mbox{ at }x) \IFF (x\in\dom(F))
\]
\end{definition}
%For every axiomatic definition, there are certain prerequisites.
A witness property for $(\varepsilon x,\varphi)$ requires an existence
hypothesis; a unique-witness property for $(\iota x,\varphi)$ requires
existence and uniqueness. A total interpretation can still assign values
when these hypotheses fail, using defaults $D_1$ and $D_2$ respectively.

A default that is also a legitimate output cannot, by its value alone,
distinguish a failed description from a successful one. For example, if
$D_1$ is the empty class, an empty output need not indicate failure. A
success predicate therefore needs information about the construction or its
context, unless the default lies outside the intended output range.
This motivates recording a validity condition alongside the value.

The relational inverse always exists as a relation. The protected-family
encoding is asymmetric, however: reversing its pairs need not produce
another protected family. Thus inversion of a representation and inversion
within a chosen representation format are different questions.

\subsection{Nonempty Values and Relational Domains}
For the next comparison, fix an injective ordered-pair encoding of sets,
and let $F\subseteq V\times V$ be a class relation. Write
\[
U_F(x)=\union\cc{y}{(x,y)\in F}
\]
for its unconditional union value, to distinguish this evaluation from the
fibre and descriptor conventions elsewhere. The domain still means the
relational domain: $x\in\Dom(F)$ if some record $(x,y)$ occurs in $F$.

\begin{definition}[Relation with nonempty evaluated values]
In this encoding, $F$ is a family of nonempty values when
\[
\forall x\in\Dom(F),\quad U_F(x)\neq\emptyclass.
\]
\end{definition}
A represented fibre may be nonempty while its evaluated value is empty.
For instance, the one-record relation $\{(x,\emptyclass)\}$ has domain
$\{x\}$ but value $\emptyclass$ at $x$. This is exactly the case that an
encoding of nonempty values alone cannot retain.

\subsection{Canonical Representations of Functions and Families}
\begin{definition}[conglomerate of all classes]
Here $\mathcal{U}$ denotes the conglomerate of all classes in the
metatheory. It is not a class-denoting term of the local theory
(MK with the proposed universe axiom).
\end{definition}
The following candidate canonicalizations encode membership using singleton
records. Their treatment of empty values distinguishes the variants:
\begin{definition}[Canonicalization of a family of nonempty classes]
\[
\canonize0(F) = \cc{(x,\{q\})}{\exists y,\ ((x,y)\in F)\land (q\in y)}
\]

\end{definition}
\begin{definition}[Canonicalization allowing empty values]
\[
\canonize1(F) = \cc{(x,\{q\})}{\exists y,\ ((x,y)\in F)\land (q\in y)}\cup\cc{(x,\emptyclass)}{x\in\Dom(F)}
\]
\[
\canonize2(F) = \cc{(x,\{q\})}{\exists y,\ ((x,y)\in F)\land (q\in y)}\cup\cc{(x,\emptyclass)}{(x\in\Dom(F))\land (F(x)=\emptyclass) }
\]
\[
\canonize3(F) = \cc{(x,s)}{\exists y,\ ((x,y)\in F)\land (s\subseteq y) \land (!q,\ q\in s)}
\]
\end{definition}

In these formulas, $F(x)$ means $U_F(x)$, and all indices and stored
components are sets. Singleton formation and replacement ensure that the
new records are sets, so class comprehension supplies the displayed
relations. Write $C_i(F)=\canonize i(F)$ for $i=0,1,2,3$.

\begin{theorem}[Values and domains of the canonicalizations]
For every set $x$ and each $i=0,1,2,3$,
\[
U_{C_i(F)}(x)=U_F(x).
\]
Moreover,
\[
\Dom(C_0(F))=\cc{x\in\Dom(F)}{U_F(x)\neq\emptyclass},
\]
whereas $\Dom(C_i(F))=\Dom(F)$ for $i=1,2,3$. In fact
$C_1(F)=C_3(F)$.
\end{theorem}
\Proof
Put $B=U_F(x)$. The nonempty records at $x$ in $C_0(F)$ are precisely
$(x,\{q\})$ for $q\in B$: membership in the union means that $q$ lies in
some original component $y$. The union of those singleton components is
$B$, and a record exists exactly when $B$ is nonempty.

The first variant adds the record $(x,\emptyclass)$ at every original
domain point. This does not change the union and restores precisely the
indices lost by $C_0$. The second adds it only when $B=\emptyclass$;
when $B\neq\emptyclass$, a singleton record already retains the index.
Thus both variants preserve the original domain and all values.

For the third variant, an at-most-one set $s\subseteq y$ is either empty
or a singleton $\{q\}$ with $q\in y$. A singleton record therefore occurs
exactly when $q\in B$, and an empty record occurs exactly when some
original component $y$ exists. These are precisely the records of $C_1(F)$.
\Done

\begin{corollary}[Idempotence and the remaining representational choice]
Each $C_i$ is idempotent: $C_i(C_i(F))=C_i(F)$. The canonicalizations
$C_1$ and $C_2$ have the same values and domain, but need not be equal.
\end{corollary}
\Proof
The records of $C_0(F)$ are determined solely by the values $U_F(x)$,
which $C_0$ preserves. The records of $C_1(F)$ and $C_2(F)$ are determined
by those values together with the domain, both of which they preserve.
This proves idempotence for the first three maps, and $C_3=C_1$ gives the
fourth case.

If $U_F(x)\neq\emptyclass$, $C_1(F)$ contains $(x,\emptyclass)$ while
$C_2(F)$ does not. For example, take $F=\{(x,\{q\})\}$ for sets $x,q$.
Both canonicalizations evaluate to $\{q\}$ at $x$, but only the first
stores the empty marker there.
\Done

The two conventions therefore solve the same value-and-domain problem
with different normal forms. Choosing one is a representation decision,
not a change to the represented class-valued family.

\begin{theorem}[Complete normal forms for values with a domain]
For class relations $F,G\subseteq V\times V$, the following are equivalent:
\begin{enumerate}
\item $\Dom(F)=\Dom(G)$ and $U_F(x)=U_G(x)$ for every set $x$;
\item $C_1(F)=C_1(G)$;
\item $C_2(F)=C_2(G)$.
\end{enumerate}
By contrast, $C_0(F)=C_0(G)$ is equivalent just to equality of all
union values, without equality of the original domains.
\end{theorem}
\Proof
Equal values determine the same singleton records. Equal domains then
determine the same empty records for $C_1$, and, together with the
values, for $C_2$. This proves the first condition implies both others.
Conversely, equality of either canonicalization implies equality of its
domain and union values, which are those of the original relation by
the preceding theorem. For $C_0$, its only records are the singletons
of members of the union values, and it preserves those values; hence
each direction of the last equivalence follows too.
\Done

Thus $C_1$ and $C_2$ each select exactly one representative of a family
specified by its domain and values. This is an assertion about the
representations, not a claim that a quotient collection of all such
class families is an object of the theory.

Here ``select'' does not require a choice function on arbitrary
collections of representatives. Each $C_i(F)$ is defined explicitly
from the records of $F$, and the theorem proves that it is independent
of which relation with the same domain and values was supplied.
Neither a well-order of the universe nor a choice of one member from
each nonempty value is used. Indeed, all members of a value contribute
singleton records. This explains why canonicalizing a family here is
a different operation from choosing a witness in the finite candidate
model of the preceding section.

\begin{theorem}[Encoding prescribed membership data]
Let $A$ be a class of set indices and $E\subseteq A\times V$ a class
relation describing the intended members of each value. Put
\[
E_x=\{q:(x,q)\in E\},\qquad
P=\{(x,\{q\}):(x,q)\in E\}
       \cup\{(x,\emptyclass):x\in A\}.
\]
Then $P$ is a protected family with domain $A$, with
$U_P(x)=E_x$ for every set $x$, and $C_1(P)=P$. It is the unique
$C_1$-fixed relation with this domain and these values.
\end{theorem}
\Proof
Every record has a set index in $A$ and an empty or singleton component.
The empty record ensures that each $x\in A$ is in the domain, and no
record has an index outside $A$. The members of the union value at $x$
are exactly the $q$ for which the singleton record $(x,\{q\})$ is
present, namely the members of $E_x$. Its records are already precisely
those prescribed by $C_1$, so $P$ is fixed. If $Q$ is another fixed
relation with the same data, the normal-form theorem gives
$Q=C_1(Q)=C_1(P)=P$.
\Done

The data $E$ are important: an arbitrary metatheoretic assignment of a
class to every index need not be definable by a class relation in the
chosen theory. The theorem assumes such a relation is supplied; it does
not assert unrestricted class-valued replacement. It does allow proper
class values. For example, for a set $a$, take $A=\{a\}$ and
$E=\{(a,q):q\in V\}$. Then $U_P(a)=V$, although every individual
record of $P$ is a set.

The next definition describes an indexed relation before a canonical
representation has been chosen.

\begin{definition}[family]
\[
(F:A\to\mathcal{U})
\IFF
(F\subseteq V\times V)
\land
(\Dom(F)=A)
\]
\end{definition}
\begin{definition}[canonical family]
One can interpret a family through $C_1$ or $C_2$, or require its
representative to be fixed by the chosen canonicalization. The theorem above
justifies preservation of values and domains, and idempotence ensures that
applying the same convention twice makes no further change.
\end{definition}



\begin{definition}[family]
In this alternative convention, a family is already a protected relation.
Its permitted second components form the class
\[
\mathsf{SS}=\{\emptyclass\}\cup\cc{\{q\}}{q\in V},
\]
consisting of the empty set and all singletons of sets. Equivalently,
these are exactly the set subsets of singletons. With this explicit
interpretation of the earlier schematic range notation $\{-\}[\,]$, put
\[
(F\family)\IFF(F\subseteq V\times\mathsf{SS}).
\]
\end{definition}
These representations have different scopes. A class relation can encode
set-indexed families with proper-class values through their members, whereas
a function symbol may also take proper classes as arguments. An ordinary
class function has set values because its graph stores ordered pairs of sets.

The singleton notation raises a related question: how should $\{V\}$ be
interpreted when proper classes cannot be elements? An empty-class default
can signal failure of singleton formation, since a successfully formed
singleton is nonempty. That reasoning does not transfer directly to function
evaluation, where the empty class may be an ordinary value. The relevant
criterion is whether the default is distinguishable from successful outputs.

\subsection{Value of a function via set--descriptors}
This subsection uses the intersection-based set descriptor defined below.
Throughout this subsection assume that $V$ is proper and that
$\intersection\emptyclass=V$, as in that construction. Its unsuccessful
value is therefore not a set, so a set-valued result certifies successful
description. This is different from the earlier union default.

\begin{definition}[Descriptor value of a relation]
Let $F$ be an arbitrary class.
\[
F(x) = \left(\iota y,\  \langle x,y\rangle\in F\right)
\labeltagd{def}{value}
\]
\end{definition}
\begin{definition}[univocal at]
\[
(F\univocalat x)\iff(!y, \langle x,y\rangle \in F)
\label{def:univocalat}
\]
\end{definition}
\begin{definition}[univocal]
\[
(F\univocal)\iff (\forall x, F\univocalat x)
\label{def:univocal}
\]
\end{definition}
\begin{theorem}[pair in function]
If $(F\univocalat x)\land(y\set)$, then
\[
\left(F(x)=y\right)\iff\left(\langle x,y\rangle\in F\right)
\label{thm:pair in function}
\]
\begin{align*}
(F\univocalat x)\land(y\set)
\implies
\left(
\left(\langle x,y\rangle\in F\right)
\implies
\left(F(x)=y\right)
\right)
\\
(y\set)
\implies
\left(
\left(F(x)=y\right)
\implies
\left(\langle x,y\rangle\in F\right)
\right)
\end{align*}
\iffalse
\[
\mbox{x set!!??}
(F\univocalat x)
\implies
\left(
\left(
\left(F(x)=y\right)\land(y\set)
\right)
\iff\left(\langle x,y\rangle\in F\right)
\right)
\]
\fi
\end{theorem}
\Proof
Write $P(z)$ for $\langle x,z\rangle\in F$, choosing $z,w$ fresh
for the fixed parameters. Use the final set-descriptor definition:
\[
Q=\{z:P(z)\land(!(w\set),P(w))\},\qquad F(x)=\intersection Q.
\]
The uniqueness guard ranges over sets; comprehension likewise retains
only set witnesses. The variable $w$ is independent of the outer
comprehension variable $z$.

First assume $P(y)$, $y\set$, and $F\univocalat x$. Local univocality
implies that at most one set satisfies $P$, and $y$ supplies such a
set witness. Thus $Q=\{y\}$. The singleton-intersection identity gives
$F(x)=\intersection\{y\}=y$, proving the forward membership-to-value
implication.

For the converse assume only $F(x)=y$ and $y\set$; no univocality
hypothesis is used. If $Q$ were empty, then $F(x)=V$, contrary to
sethood of $y$ and properness of $V$. Choose $a\in Q$. Comprehension
gives $a\set$, $P(a)$, and at most one set satisfying $P$. Every
member of $Q$ is such a set, so $Q=\{a\}$. Hence
$y=F(x)=\intersection Q=a$, and equality substitution gives $P(y)$,
which is $\langle x,y\rangle\in F$. This proves the separately
stated reverse implication with only its sethood hypothesis.
\Done

\begin{theorem}[Witness property of a defined set descriptor]
\[
((\iota y,\ \phi)\ldefined)
\implies
\phi[(\iota y,\ \phi)/y]
\]
\end{theorem}
\Proof
Use the local convention that a set descriptor is defined exactly when its
value is a set. Put
\[
Q=\cc{y}{\phi\land(!(y\set),\phi)}.
\]
Its intersection is the value of the descriptor. If $Q$ were empty, this
value would be $V$, which is a proper class. Definedness therefore implies
that $Q$ is nonempty. Its uniqueness condition makes $Q$ a singleton,
say $Q=\{a\}$. Then $\intersection Q=a$, and membership of $a$ in $Q$
gives the required instance of $\phi$. Substitution is understood to avoid
capture of free variables.
\Done

\begin{remark}[Only uniqueness among sets is needed]
\label{thm:pair in function alternative}
In the membership-to-value implication, the stronger assumption
$F\univocalat x$ can be replaced by
$!(w\set),\ \langle x,w\rangle\in F$. The proof uses exactly this
set-restricted uniqueness together with the given set witness $y$.
Even if the predicate also holds of a proper class, that additional
object is not collected by $Q$ and does not affect the descriptor.
The reverse implication instead uses sethood of the result to exclude
$V$; it needs neither uniqueness assumption in advance and would not
be valid with an arbitrary set default.
\end{remark}

\subsection{Earlier Set-Descriptor Axioms}
\begin{axiom}[Choice witness]
\[
\exists (x\set),\ \varphi
\implies
(\varepsilon x,\  \varphi) \in \cc{x}{\varphi}
\]
\end{axiom}
\begin{axiom}[Sethood signals existence]
\[
(\varepsilon x,\  \varphi) \set
\implies
\exists (x\set),\ \varphi
\]
\end{axiom}

\begin{axiom}[Unique description witness]
\[
\exists! (x\set),\ \varphi
\implies
(\iota x,\  \varphi) \in \cc{x}{\varphi}
\]
\end{axiom}
\begin{axiom}[Sethood signals unique existence]
%(\iota x,\  \varphi) \in \cc{x}{\varphi}
\[
(\iota x,\  \varphi) \set
\implies
\exists! (x\set),\ \varphi
\]
\end{axiom}



The next construction asks whether an intersection can describe a unique
class, including a proper class:
\[
(I y,\varphi)=\intersection\cc{y}{\varphi\land(!y,\varphi)}.
\]
The obstacle is local to this encoding: class comprehension collects only
set witnesses. If the unique satisfying object is a proper class, it is not
collected. Moreover, uniqueness among all classes and uniqueness among sets
are different conditions when $\varphi$ has witnesses of both kinds.
The construction below therefore restricts uniqueness to set witnesses.
A descriptor for arbitrary classes would require a different construction
or an explicitly enlarged semantic setting.
\begin{definition}[iota operator for sets]
\[
(\iota y,\  \varphi) = \intersection \cc{y}{\varphi \land (!(y\set),\ \varphi)}
\]
\label{def:iota operator}
\end{definition}

\begin{theorem}[Sethood characterizes unique set description]

%Note: it should be possible to substitute some fresh variable $w$ instead of $\varphi$.
%This may not be true:
%\[ ???
%((\iota y,\  \varphi) \set)
%\iff
%(\exists!(b\set),\ \varphi)
%\]
%Variant2(true before redef but ugly):
%\[
%((\iota y,\  \varphi) \set)
%\iff
%(\exists(b\set),\ \varphi)
%\land
%(!b,\ \varphi)
%\]
%Variant3:
%\[
%((\iota y,\  \varphi) \set)
%\iff
%(\exists!b,\ \varphi)
%\]
%Variant4:
%\[
%((\iota (y\set),\  \varphi) \set)
%\iff
%(\exists!(y\set),\ \varphi)
%\]
%Variant5:
\[
((\iota y,\  \varphi) \set)
\iff
(\exists!(y\set),\ \varphi)
\]
\label{thm:iota set exists unique set}
\end{theorem}
\Proof
Put $Q=\cc{y}{\varphi\land(!(y\set),\varphi)}$. By definition,
$\iota y,\varphi=\intersection Q$.

If there is exactly one set witness $a$, the uniqueness condition holds and
$Q=\{a\}$. The singleton-intersection identity gives
$\intersection Q=a$, which is a set.

Conversely, suppose $\intersection Q$ is a set. Since
$\intersection\emptyclass=V$ and $V$ is proper, $Q$ cannot be empty.
Choose $a\in Q$. By comprehension, $a$ is a set satisfying $\varphi$, and
the second conjunct in the definition of $Q$ states that there is at most
one such set. Thus there is exactly one set witness.
\Done

\begin{corollary}[The successful value and the default]
If $a$ is the unique set satisfying $\varphi$, then
$(\iota y,\varphi)=a$. If there is no unique set satisfying $\varphi$,
then $(\iota y,\varphi)=V$.
\end{corollary}
\Proof
The successful case was proved above. In the other case, a member of $Q$
would supply both existence and uniqueness, a contradiction. Hence
$Q=\emptyclass$ and its intersection is $V$.
\Done

For example, a predicate with the empty set as its unique set witness
returns $\emptyclass$, whereas a predicate with no set witness returns $V$.
If a predicate holds of exactly one set and also of a proper class, this
set-restricted descriptor still returns the set. This example explains why
its uniqueness quantifier must explicitly range over sets: it solves the
set-description problem, not the different problem of selecting a unique
class witness.

\ifdefined\qwertyflag
\undef\qwertyflag
\end{document}
\fi
